(1) Jacobi–Anger公式 \(e^{ikr\cos\varphi}=\sum_{m\in\mathbb Z}i^mJ_m(kr)e^{im\varphi}\)。取
\(p=p_{\rm in}+\sum_m i^mb_mH_m^{(1)}(kr)e^{im\varphi}\)，在圆柱 \(r=a\) 处对 \(r\) 微分，每个角向系数独立为零，故
\[
\boxed{b_m=-\frac{J_m'(ka)}{H_m^{(1)\prime}(ka)}}.
\]
这里撇号对无量纲变量求导，额外的 \(k\) 因子在分子分母抵消；刚壁采用Neumann而非软壁Dirichlet条件。

(2) 当 \(kr\gg1\)，有 \(H_m^{(1)}(kr)\sim\sqrt{2/(\pi kr)}\,e^{i(kr-m\pi/2-\pi/4)}\)。故远场各模散射角振幅的角傅里叶系数是 \(b_m\)，角积分的正交性给每单位轴长总散射截面
\[
\boxed{\sigma_{\rm sc}=\frac4k\sum_{m=-\infty}^{\infty}|b_m|^2}.
\]
由于圆柱中心对称且 \(J_{-m}=(-1)^mJ_m,\ H_{-m}=(-1)^mH_m\)，有 \(b_{-m}=b_m\)，于是 \(m>0\) 每项出现两次，不可误除简并系数。

(3) 写 \(z=ka\ll1\)。\(J_0'\sim-z/2\)、\(H_0^{(1)\prime}\sim2i/(\pi z)\)，所以 \(b_0\sim-i\pi z^2/4\)。又 \(J_1'\sim1/2\)、\(H_1^{(1)\prime}\sim2i/(\pi z^2)\)，故 \(b_{\pm1}\sim i\pi z^2/4\)。更高 \(m\) 从更高幂 \(z^{2|m|}\) 起。因此三个低阶模都以 \(z^4\) 贡献强度，
\[
\boxed{\sigma_{\rm sc}\sim \frac4k\frac{3\pi^2z^4}{16}
=\frac{3\pi^2}{4}k^3a^4}.
\]
单取 \(m=0\) 会漏掉 \(2/3\) 的主导总量，这与压缩性不同的软声学圆柱低频散射不应混同。

(4) 因 \(J_m,Y_m\) 在实 \(z>0\) 为实函数，\(H_m^{(1)\prime}=J_m'+iY_m'\)，由第一问 \(S_m=1+2b_m=(H_m'-2J_m')/H_m'=(-J_m'+iY_m')/(J_m'+iY_m')=-H_m^{(2)\prime}/H_m^{(1)\prime}\)。二者互为复共轭，所以 \(\boxed{|S_m|^2=1}\)。这表示每一个部分波的入射径向能流被完全弹性散射到出射通道，而无能量吸收；位相可变但模不变。
